% Autorska priprava k Realisticky.cz 020201; ulohy jsou autorske a obmenene.
% Zamerne nenahrazuje vyklad presnymi kopiemi cizich textu/prikladu.
\documentclass[11pt,a4paper]{article}
\input{../../sablony/pripravy/priprava.sty}
\usepackage{tikz}
\nastavPripravu{1. ročník}{Rovnice a nerovnice}{72}{Rovnice a ekvivalentní úpravy}
\begin{document}
\blok{Návaznost a motivace}
Z lineárních funkcí známe otázku, kdy mají dvě funkce stejnou hodnotu. Rovnice $L(x)=P(x)$ hledá právě takové přípustné argumenty. V algebře už žáci přičítali stejné výrazy na obou stranách identity; nyní to zdůvodníme při hledání kořenů.
\begin{definiceBox}
\textbf{Kořen rovnice}. Je číslo z určeného oboru, pro které se levá a pravá strana rovnice rovnají. Množinu všech kořenů označujeme $K$.
\end{definiceBox}
\begin{definiceBox}
\textbf{Ekvivalentní rovnice}. Dvě rovnice na stejném oboru jsou ekvivalentní, pokud mají stejnou množinu řešení. Symbol $\Leftrightarrow$ značí, že příslušný krok zachovává řešení.
\end{definiceBox}
\blok{Princip a odvození}
Na obou stranách lze přičíst nebo odečíst stejný výraz definovaný pro všechna uvažovaná $x$; násobení a dělení je ekvivalentní jen pro nenulový činitel. Podmínky není dovoleno ztratit. V zápisu $\Leftrightarrow$ stojí matematické zdůvodnění, nikoli dekorace.
\blok{Řešené příklady: od jednoduchého ke složitějšímu}
\begin{prikladReseniBox}[title={Příklad 1}]
Vyřešte $3x-5=10$ a proveďte zkoušku.
\tcblower
\begin{align*}
3x-5 &= 10 \\
 3x &= 15 \\
 x &= 5
\end{align*}

Zkouška: $3\cdot5-5=10$, tedy $K=\{5\}$. Ekvivalentní byly oba kroky.
\end{prikladReseniBox}
\begin{prikladReseniBox}[title={Příklad 2}]
Vyřešte $\frac{2x-1}{3}=\frac{x+4}{2}$.
\tcblower
\begin{align*}
2(2x-1) &= 3(x+4) \\
 4x-2 &= 3x+12 \\
 x &= 14
\end{align*}

Oba původní jmenovatele jsou nenulové konstanty; kontrola: $27/3=18/2=9$.
\end{prikladReseniBox}
\poznamka{Zdůrazni změnu celé levé i celé pravé strany. Požaduj vždy jednu zkoušku dosazením do původní rovnice.}
\blok{Samostatné procvičování}
\begin{zadaniReseniBox}[title={Úloha 1}]
Vyřešte $x+7=3$.
\tcblower
\begin{align*}
x+7 &= 3 \\
 x &= 3-7 \\
 x &= -4
\end{align*}
\end{zadaniReseniBox}
\begin{zadaniReseniBox}[title={Úloha 2}]
Vyřešte $4(x-1)=2x+6$.
\tcblower
\begin{align*}
4x-4 &= 2x+6 \\
 2x &= 10 \\
 x &= 5
\end{align*}
\end{zadaniReseniBox}
\begin{zadaniReseniBox}[title={Úloha 3}]
Posuďte $5x-4=5x+3$.
\tcblower
\begin{align*}
5x-4 &=5x+3 \\
 -4 &=3 \\
 K&=\varnothing
\end{align*}
\end{zadaniReseniBox}
\blok{Volitelná rezerva}
\begin{zadaniReseniBox}[title={Úloha 1}]
Rozhodněte, zda je násobení rovnice $x=1$ výrazem $x-1$ ekvivalentní.
\tcblower
\begin{align*}
x &= 1 \\
 x(x-1) &= x-1 \\
 (x-1)^2 &=0
\end{align*}

Bez podmínky na násobící výraz tuto úpravu nesmíme označit za ekvivalentní pro všechny rovnice; zde konkrétně kořen náhodou nezmění.
\end{zadaniReseniBox}
\blok{Domácí úkol}
\begin{zadaniReseniBox}[title={Domácí příprava}]
\begin{enumerate}[label=\alph*)]
\item
$\frac{x}{3}+2=5$.
\item
$2(x+3)=2x+6$.
\item
$7x-2=3x+10$.
\item
$5x-9=16$.
\item
$4(x-2)=2x+6$.
\item
$3(2x-1)=6x+4$.
\end{enumerate}
\tcblower
\begin{enumerate}[label=\alph*)]
\item
\begin{align*}
\frac{x}{3} &=3 \\
 x &=9 \\
 K&=\{9\}
\end{align*}
\item
\begin{align*}
2x+6 &=2x+6 \\
 0 &=0 \\
 K&=\mathbb R
\end{align*}
\item
\begin{align*}
4x &=12 \\
 x &=3 \\
 K&=\{3\}
\end{align*}
\item
$\begin{aligned}
5x-9&=16\\5x&=25\\x&=5
\end{aligned}$
\item
$\begin{aligned}
4x-8&=2x+6\\2x&=14\\x&=7
\end{aligned}$
\item
$\begin{aligned}
6x-3&=6x+4\\-3&=4\quad\text{-- spor}\\K&=\varnothing
\end{aligned}$
\end{enumerate}
\end{zadaniReseniBox}
\poznamka{Orientační časové rozvržení: 5 + 12 + 12 + 10 + 6 = 45 min. Přednost má pochopení jádra; rezerva není součástí časového součtu. Domácí úlohy jsou určeny na samostatnou práci, řešení zůstávají učiteli.}
\end{document}
